% Source correction, 2 October 2026. The superseded v39 packet is preserved.
\begin{definition}
\label{def:p031-ordinary-locally-split}
Let $\pi:E\to X$ be a finite surjective morphism of $K$-analytic spaces.
We call $\pi$ \emph{ordinary-locally split} if every $x\in X$ has an
ordinary open neighborhood $U$ for which
\begin{equation}
\label{eq:p031-ordinary-local-splitting}
 \pi^{-1}(U)=\coprod_{i=1}^{d}U_i,
 \qquad \pi_i:=\pi|_{U_i}:U_i\xrightarrow{\ \sim\ }U
\end{equation}
as $K$-analytic spaces. The fiber cardinality is then locally constant.
\end{definition}

\begin{remark}[Real coefficients and finite maxima]
\label{thm:p031-lfs-fiber-maximum}
For a real-valued continuous function $u$ on $E$, the function
\begin{equation}
\label{eq:p031-lfs-fiber-maximum}
 \pi_{\max}u(x):=\max_{e\in\pi^{-1}(x)}u(e)
\end{equation}
is continuous in this splitting range. If $\sigma_i:=\pi_i^{-1}$, then
\begin{equation}
\label{eq:p031-local-fiber-maximum}
 \pi_{\max}u|_U=\max_i\{u|_{U_i}\circ\sigma_i\}.
\end{equation}
It need not be real PL, even when $u$ is LFS. Indeed,
\cite[Example~5.3]{GJR26} gives $h=\pi\log|t|+1$ on a one-dimensional
torus over a field with additive value group $\mathbf Q$, with
$\max\{h,0\}$ not real PL. Here $h$ is LFS: if $-\log|a|=1$, then
$h=\pi(-\log|t^{-1}|)-\log|a|$ is a nonnegative real linear combination
of unit-logarithm functions, as in \cite[\S7.9]{GJR26}. Apply the fold map
$X\amalg X\to X$ to the function equal to $h$ and $0$ on its two sheets.
Thus the real-coefficient description in \S7.9 does not imply closure
under arbitrary finite maxima. The rational-coefficient closure used
below has a separate common-denominator proof.
\end{remark}

\begin{theorem}[LFS envelopes on a splitting neighborhood]
\label{thm:p031-lfs-envelope-pullback}
Let $\pi:E\to X$ be ordinary-locally split between good strict
$K$-analytic spaces, fix a splitting neighborhood $U$, and put
$E_U:=\pi^{-1}(U)$. For any $f:U\to\mathbf R\cup\{-\infty\}$, let
$s_{\mathrm{LFS},U}$ and $s_{\mathrm{LFS},E_U}$ be the raw envelopes
of LFS minorants defined on all of $U$ and $E_U$, respectively. Then
\begin{equation}
\label{eq:p031-lfs-envelope-pullback}
 s_{\mathrm{LFS},E_U}=s_{\mathrm{LFS},U}\circ\pi.
\end{equation}
The two families are empty simultaneously, with empty supremum $-\infty$.
This compares envelopes formed on the displayed domains. It does not
identify restrictions of envelopes formed globally on $X$ and $E$.
\end{theorem}
